\chapter{Problem-Solving Patterns and Classical Paradoxes}
\label{ch:patterns}

This closing chapter collects the recurring \emph{techniques} for solving probability problems and
the famous paradoxes that test them. Each paradox is resolved by identifying exactly which assumption
was mishandled, which is the most reliable way to build intuition.

\section{A toolkit of patterns}

\begin{enumerate}
\item \textbf{Count the complement.} ``At least one'' is $1-\Prob(\text{none})$. Birthday problem:
$1-\prod_{k=0}^{n-1}(1-k/365)\approx1-e^{-n(n-1)/730}$, exceeding $\tfrac12$ at $n=23$.
\item \textbf{Linearity of expectation with indicators.} Break the count into indicators; dependence is
irrelevant. Expected fixed points of a random permutation is $n\cdot\frac1n=1$ for all $n$; expected
number of distinct values in $n$ draws from $m$ equally likely is $m\bigl(1-(1-\tfrac1m)^n\bigr)$.
\item \textbf{Symmetry.} If $X,Y$ are exchangeable, $\Prob(X>Y)=\tfrac12(1-\Prob(X=Y))$ with no computation.
The probability that the $k$th card dealt is an ace equals that of the first.
\item \textbf{Condition on the first step} (first-step analysis, Chapter~\ref{ch:condexp}). Expected flips for
HH: $E=\tfrac12(1+E)+\tfrac14(2+E)+\tfrac14\cdot2\Rightarrow E=6$.
\item \textbf{Condition on the right variable.} Compute $\Exp[Y]=\Exp[\Exp[Y\given X]]$ with $X$ chosen to make
the inner expectation easy.
\item \textbf{Reflection and bijections.} Convert hard counts into easy ones by a one-to-one
map (ballot problem, Chapter~\ref{ch:markov}).
\item \textbf{Martingales.} When a quantity ``has no drift'', optional stopping gives answers in one line.
\item \textbf{Generating functions.} Convert convolutions into products and recurrences into algebra.
\item \textbf{Bounds when exact is hard.} Union bound, Markov/Chebyshev, Chernoff; the probabilistic method:
if $\Exp[X]\ge c$ then some outcome has $X\ge c$.
\item \textbf{Change of order.} $\Exp X=\int\Prob(X>x)dx$; swap sum and integral (Tonelli).
\item \textbf{Sanity checks.} Small cases ($n=1,2$), limits ($p\to0,1$), units/symmetry, and the sum of
probabilities equal to $1$.
\end{enumerate}

\begin{interview}{a standard solving routine}
(1) State the sample space and what ``uniformly random'' means. (2) Identify an event or random variable.
(3) Pick a pattern above. (4) Verify with a small case. Many ``paradoxes'' are really failures of step (1).
\end{interview}

\subsection*{Three worked problems}

\begin{example}[Expected maximum of $n$ uniforms]
With $M=\max U_i$, $\Prob(M\le x)=x^n$, so $\Exp M=\int_0^1(1-x^n)dx=1-\frac1{n+1}=\frac n{n+1}$.
\end{example}

\begin{example}[Broken stick triangle]
Break a unit stick at two independent uniform points. A triangle forms iff every piece is $<\tfrac12$.
By inclusion--exclusion over the events $\{$piece $i\ge\tfrac12\}$, which are mutually exclusive (two
pieces cannot both be $\ge\tfrac12$ unless they are exactly $\tfrac12$), each with probability $\tfrac14$
(the vector of pieces is uniform on the simplex, and $\Prob(\text{piece}_1\ge\tfrac12)=(1-\tfrac12)^2$), the
probability is $1-3\cdot\tfrac14=\tfrac14$.
\end{example}

\begin{example}[Random walk returns]
Probability that a symmetric walk on $\Z^d$ ever returns to the origin is $1$ for $d=1,2$ and
strictly less than $1$ for $d\ge3$ (P\'olya; about $0.34$ in $d=3$). The reason: the return probability at step
$2n$ is $\asymp n^{-d/2}$ and $\sum n^{-d/2}$ converges iff $d\ge3$ (Chapter~\ref{ch:markov}). ``A drunk man
finds his way home, a drunk bird may be lost forever.''
\end{example}

\section{Classical paradoxes}

\subsection{Monty Hall}

Three doors, one prize. You pick door 1. The host, who \emph{knows} where the prize is, always opens a
different door showing a goat, then offers a switch. By Bayes with prize position $P$ uniform and the host
choosing uniformly between two goat doors when he has a choice:
\[
\Prob(P=2\given\text{host opens }3)=\frac{\tfrac13\cdot1}{\tfrac13\cdot\tfrac12+\tfrac13\cdot1+\tfrac13\cdot0}=\frac23 .
\]
Switching wins with probability $\tfrac23$.

\begin{pitfall}{the host's protocol is part of the problem}
If the host opens a random door and it \emph{happens} to show a goat, switching gives only $\tfrac12$.
The data ``door 3 has a goat'' is the same; the likelihood of the data differs. This is the general
moral: always model the \emph{process that generated the information}.
\end{pitfall}

\subsection{Simpson's paradox}

Treatment A can have a higher success rate than B within each of two subgroups, yet a lower rate overall.
Example: A: $\tfrac{81}{87}$ (mild) and $\tfrac{192}{263}$ (severe) versus B: $\tfrac{234}{270}$ and $\tfrac{55}{80}$.
Within each group A wins ($93.1\%>86.7\%$, $73.0\%>68.8\%$), but overall $\tfrac{273}{350}=78\%<\tfrac{289}{350}=83\%$.
The reason is arithmetic: the overall rate is a weighted average with \emph{different weights}
(A treats mostly the hard cases). Formally $\Prob(S\given T)=\sum_g\Prob(S\given T,g)\Prob(g\given T)$ and the
weights $\Prob(g\given T)$ depend on $T$. Which comparison is correct depends on the causal structure: if the group
is a confounder (severity influenced the choice of treatment), use the stratified comparison.

\subsection{Bertrand's paradox}

``Pick a random chord of a circle; what is the probability it is longer than the side of the inscribed
equilateral triangle?'' Three natural definitions of ``random'' give three answers:
\begin{itemize}
\item random endpoints on the circle: $\tfrac13$;
\item random radius, then a random point on it as the midpoint: $\tfrac12$;
\item random midpoint uniform in the disc: $\tfrac14$.
\end{itemize}
\textit{Resolution}: ``random chord'' does not specify a probability space. The three methods define
different measures on the set of chords. Only with an invariance requirement (rotation, translation and
scale invariance, as in Jaynes' argument) does the second answer, $\tfrac12$, become singled out.

\subsection{The two-envelope paradox}

Two envelopes contain $x$ and $2x$. You hold one with amount $A$ and reason: the other holds $2A$ or $A/2$
with probability $\tfrac12$ each, so its expectation is $\tfrac54A>A$ and you should always switch, an
absurdity (by symmetry switching back is again profitable).

\textit{Resolution}: Let the smaller amount be $X$. The step ``the other is $2A$ or $A/2$ with probability $\tfrac12$
each \emph{conditional on} $A=a$'' implies a prior on $X$ with $\Prob(X=a)=\Prob(X=a/2)$-type
consistency for every $a$, i.e.\ a uniform ``improper'' prior on the powers of $2$, which cannot be a probability
distribution. For any proper prior, $\Exp[B\given A=a]$ exceeds $a$ only for small $a$ and falls below it for
large $a$. And if $\Exp[X]<\infty$, unconditional expected gain from switching is exactly $0$. If $\Exp[X]=\infty$
(a valid distribution), expectations of $A$ and $B$ are both infinite and the comparison is meaningless.

\subsection{St.\ Petersburg paradox}

Toss a fair coin until the first head, at toss $N$; win $2^N$. Then
$\Exp[\text{win}]=\sum_{n\ge1}2^{-n}2^n=\sum1=\infty$, yet nobody would pay much to play. Resolutions:
\begin{itemize}
\item \emph{Finite bankroll}: if the casino can pay at most $2^k$, the value is $\approx k/2+1$ (log of the cap).
\item \emph{Utility}: with $u(w)=\ln w$ (Bernoulli), $\Exp\ln(\text{win})=\sum2^{-n}n\ln2=2\ln2$, finite, certainty
equivalent $4$.
\item \emph{Typical behaviour}: the average of $n$ plays grows like $\log_2n$ (a weak law with truncation---Feller),
never settling to a limit.
\end{itemize}
The paradox warns that expectation is the wrong summary when the distribution has an infinite mean.

\subsection{The boy--girl problems}

``A family has two children; at least one is a boy. Probability both are boys?'' The sample space
$\{BB,BG,GB,GG\}$ restricted to ``at least one B'' gives $\tfrac13$. But if the information is ``the elder is a boy''
the answer is $\tfrac12$; if ``I met one child, a boy'' (random child), it is $\tfrac12$. The wording encodes
\emph{how the information was selected}, once more the Monty Hall lesson.

\subsection{Sleeping Beauty and the prosecutor's fallacy}

The \textbf{prosecutor's fallacy} confuses $\Prob(\text{evidence}\given\text{innocent})$ with
$\Prob(\text{innocent}\given\text{evidence})$. If a DNA match has false-match probability $10^{-6}$ in a
city of $10^7$, about $10$ innocent matches are expected; given only the match, the chance that a matching person
is the culprit is about $\tfrac1{11}$, not $1-10^{-6}$. Always apply Bayes with the prior.

\subsection{Non-transitive dice and the Condorcet cycle}

Dice A: $\{2,2,4,4,9,9\}$, B: $\{1,1,6,6,8,8\}$, C: $\{3,3,5,5,7,7\}$. Enumerating the $9$ equally likely pairs
of faces (each face value is doubled, so the distribution is unaffected) gives $\Prob(A>B)=\tfrac59$,
$\Prob(B>C)=\tfrac59$, $\Prob(C>A)=\tfrac59$. ``Beats'' is not transitive, so no die is best: the relation
$\Prob(X>Y)>\tfrac12$ is not an order on random variables, unlike $\Exp X>\Exp Y$.

\subsection{Friendship and regression to the mean}

The \textbf{friendship paradox}: your friends have on average more friends than you, because a friend is sampled
proportionally to degree (size-biasing, cf.\ the inspection paradox, Chapter~\ref{ch:poisson}): the expected degree of a
random friend is $\Exp[D^2]/\Exp[D]\ge\Exp[D]$. \textbf{Regression to the mean}: if $X,Y$ are standardised with
correlation $\rho<1$, then $\Exp[Y\given X=x]=\rho x$, closer to the mean than $x$; extreme first performances are
followed by less extreme ones with no causal intervention (Chapter~\ref{ch:jointcont}).

\section{Practice list for interviews}

Each should be solvable in minutes with the toolkit above.
\begin{enumerate}
\item Expected number of rolls of a fair die to see all six faces ($14.7$; coupon collector).
\item Probability a random permutation of $n$ has no fixed point ($\to e^{-1}$).
\item Expected length of the longest run of heads in $n$ tosses ($\approx\log_2n$).
\item Two players alternate drawing balls without replacement; who is favoured? (Each position equally likely
by symmetry.)
\item A stick is broken at a random point; expected ratio of the shorter to the longer ($2\ln2-1\approx0.386$).
\item Choose $n$ points on a circle; probability they lie on a semicircle ($n/2^{n-1}$).
\item From a stream of unknown length, keep one uniformly random item with $O(1)$ memory (reservoir sampling: replace the
current item with the $k$th with probability $1/k$).
\item Gambler starts with $\$k$, wins/loses $\$1$ fairly until $0$ or $N$; ruin probability $1-k/N$.
\item Expected time for a fair coin to show three heads in a row ($14$).
\item The secretary problem: stop at the best after rejecting the first $n/e$.
\end{enumerate}

\section*{Exercises}
\addcontentsline{toc}{section}{Exercises}

\begin{enumerate}
\item Derive $n/2^{n-1}$ for points on a semicircle by conditioning on the point that is the ``first'' in clockwise
order of the semicircle.
\item Prove the reservoir sampling claim by induction.
\item Verify the Simpson's paradox numbers above.
\item Show Bertrand's three probabilities by direct integration.
\item Show $\Exp[D^2]\ge(\Exp D)^2$ and interpret the friendship paradox.
\item Find three non-transitive dice with all faces from $\{1,\dots,9\}$ other than those above, and compute their cycle
probability.
\item (Challenge) Show that in the two-envelope setting with a geometric prior on the smaller amount, switching is
profitable only for small observed amounts.
\end{enumerate}
